\documentclass[12pt]{article}
\usepackage[utf8]{inputenc}
\usepackage[T1]{fontenc}
\usepackage{hyperref}
\title{Tlachichilco Tepehua Language Resources for K-12 Curriculum Development}
\author{Compiled by Claude Sonnet 5 (claude-sonnet-5) for Andr\'as Kornai}
\date{}
\begin{document}
\maketitle

\section{Introduction}
Tlachichilco Tepehua (ISO 639-3: tpt; Glottocode tlac1235) is a Totonacan language of Veracruz, Mexico, with about 3,000--4,700 speakers depending on the source. It is endangered. Search found scripture and audio material and a few reference entries. It found no grammar, dictionary or school materials.

\section{Primary Educational Resources}
None found.

\section{Linguistic and Reference Materials}
\begin{itemize}
\item PHOIBLE phonological inventory (34 segments): \url{https://phoible.org/languages/tpt}
\item Glottolog entry: \url{https://glottolog.org/resource/languoid/id/tlac1235}
\item Verbix endangered-languages entry: \url{https://docs.verbix.com/EndangeredLanguages/TlachichilcoTepehua}
\item Open-access article on pronunciation of Tepehua as a second language in Tlachichilco (in Spanish, about Spanish-first speakers learning Tepehua): \url{https://revistas.ucm.es/index.php/CLAC/article/view/88219}
\item Endangered Languages Project resources page: \url{https://www.endangeredlanguages.com/node/38943/resources}
\end{itemize}

\section{Religious and Community Resources}
\begin{itemize}
\item ScriptureEarth: New Testament in Tlachichilco Tepehua with text, audio and PDF: \url{https://www.scriptureearth.org/00eng.php?iso=tpt}
\item ScriptureEarth index page: \url{https://www.scriptureearth.org/00i-Scripture_Index.php?iso=tpt}
\end{itemize}
No community organisations were found.

\section{Digital and Audio Resources}
\begin{itemize}
\item Global Recordings Network audio: \url{https://globalrecordings.net/en/language/tpt}
\end{itemize}

\section{Curriculum Development Context}
The scripture text and audio are the main written and spoken material available. A curriculum would require new documentation and materials developed with the community.
\end{document}